\documentclass{article}

% if you need to pass options to natbib, use, e.g.:
%     \PassOptionsToPackage{numbers, compress}{natbib}
% before loading aaj2026


% ready for submission
\usepackage[most]{tcolorbox}
\usepackage[final]{aaj2026}
\bibliographystyle{plainnat}
\usepackage[numbers]{natbib}

% to compile a preprint version, e.g., for submission to arXiv, add add the
% [preprint] option:
%     \usepackage[preprint]{aaj2026}


% to compile a camera-ready version, add the [final] option, e.g.:
%     \usepackage[final]{aaj2026}


% to avoid loading the natbib package, add option nonatbib:
%    \usepackage[nonatbib]{aaj2026}


\usepackage[utf8]{inputenc} % allow utf-8 input
\usepackage[T1]{fontenc}    % use 8-bit T1 fonts
\usepackage{hyperref}       % hyperlinks
\usepackage{url}            % simple URL typesetting
\usepackage{booktabs}       % professional-quality tables
\usepackage{amsfonts}       % blackboard math symbols
\usepackage{nicefrac}       % compact symbols for 1/2, etc.
\usepackage{microtype}      % microtypography
\usepackage{xcolor}         % colors


\title{Born Displaced: Survival, Conflict and the Children Left Behind}


% The \author macro works with any number of authors. There are two commands
% used to separate the names and addresses of multiple authors: \And and \AND.
%
% Using \And between authors leaves it to LaTeX to determine where to break the
% lines. Using \AND forces a line break at that point. So, if LaTeX puts 3 of 4
% authors names on the first line, and the last on the second line, try using
% \AND instead of \And before the third author name.


\author{%
  Gabriele Loreti \\
  Department of AI and Decision Sciences\\
  Luiss Guido Carli\\
  \texttt{gabriele.loreti@studenti.luiss.it} \\
  \And
  Matteo Perrucci \\
  Department of AI and Decision Sciences\\
  Luiss Guido Carli\\
  \texttt{matteo.perrucci@studenti.luiss.it} \\
  \AND
  Giulio Presaghi \\
  Department of AI and Decision Sciences\\
  Luiss Guido Carli\\
  \texttt{giulio.presaghi@studenti.luiss.it} \\
  \And
  Chiara Valentini \\
  Department of AI and Decision Sciences\\
  Luiss Guido Carli\\
  \texttt{chiara.valentini@studenti.luiss.it} \\
}
  % \AND
  % Coauthor \\
  % Affiliation \\
  % Address \\
  % \texttt{email} \\
  % \And
  % Coauthor \\
  % Affiliation \\
  % Address \\
  % \texttt{email} \\
  % \And
  % Coauthor \\
  % Affiliation \\
  % Address \\
  % \texttt{email} \\
}


\begin{document}


\maketitle


\begin{abstract}
This study examines the systematic relationship between forced displacement and child mortality across a global panel of countries from 2010 to 2022. Amidst an era of unprecedented humanitarian crises, we investigate whether the upheaval of populations directly correlates with increased child mortality rates and to what extent macroeconomic factors, specifically GDP per capita, unemployment, and inflation, mediate this link. Using a longitudinal fixed-effects framework, our analysis isolates the health outcomes of displacement while accounting for economic destabilization as a primary causal pathway. Preliminary findings suggest that while forced displacement significantly elevates mortality risks, the severity of this impact is heavily conditioned by a nation’s economic resilience. These results highlight the critical role of macroeconomic stabilization in mitigating the public health consequences of humanitarian displacement.
\end{abstract}

\section{Introduction and Question}

This report presents the work of a four-person data analytics team---Chiara Valentini, Giulio Presaghi, Matteo Perrucci, and Gabriele Loreti---commissioned by the IFRC Data \& Analytics Unit to support the humanitarian campaign \textit{Born Displaced}, aimed at raising funds for pediatric hospitals in countries affected by forced displacement and high child mortality. Our role was to identify, through data analysis and visualization, which countries most urgently require intervention, and to produce a set of interactive visualizations for embedding in the IFRC campaign webpage.

The central research question is twofold: does a systematic relationship exist between forced displacement and child mortality (ages 1 to 14) across countries in the period 2010 to 2022, and do macroeconomic conditions---specifically GDP per capita, unemployment rate, and inflation---mediate this relationship?

The urgency is grounded in documented evidence. By end of 2022, 108.4 million people were forcibly displaced worldwide, the largest figure ever recorded (UNHCR, 2022). Children represent 41\% of all displaced people despite making up only 30\% of the global population, and those in conflict-affected settings are almost three times more likely to die before age five than children elsewhere (UNICEF, 2022). The temporal window of 2010 to 2022 was chosen to capture four major crises---the Syrian civil war, the collapse of South Sudan, the Venezuelan exodus, and the 2022 Ukrainian displacement surge---while maintaining data comparability across sources.

The system serves two audiences: IFRC communication and advocacy officers who need data-driven evidence to design campaigns and allocate resources, and the general public who need accessible and emotionally resonant visualizations to understand the crisis and act. It supports five analytical tasks:

\begin{itemize}[leftmargin=2em,itemsep=2pt]
  \item \textbf{T1}: identifying countries with the highest combined displacement and child mortality burden at any given year;
  \item \textbf{T2}: tracking the displacement--mortality relationship over 2010 to 2022;
  \item \textbf{T3}: exploring the individual crisis profile of a specific country across displacement, mortality, and GDP over time;
  \item \textbf{T4}: identifying countries most in need of intervention based on a composite Crisis Index for 2019 to 2022;
  \item \textbf{T5}: understanding how country crisis profiles changed over time, including sudden shocks such as the 2022 Ukraine invasion.
\end{itemize}

\section{Data}

The analysis draws on three publicly available datasets: the UNHCR Refugee Data Finder (annual country-level counts of refugees, asylum seekers, IDPs, stateless persons, and returned refugees, 2010 to 2022), the UNICEF Child Mortality Database (annual mortality rate estimates for ages 1 to 4 and 5 to 14 with confidence bounds, 195 countries), and the World Bank Open Data platform (GDP per capita, unemployment rate, and inflation CPI, 2010 to 2022).

All three were integrated using Python (\texttt{pandas}), merged on ISO3 country code and year. The UNICEF dataset required parsing of compound geographic identifiers. The final unified dataset contains 2{,}600 rows across 19 columns covering 195 countries over 13 years.

Key cleaning decisions included removing 42 regional aggregates from the UNICEF dataset, filling \texttt{total\_displaced} with zero for 23 small states where displacement is structurally absent, leaving South Sudan 2010 and Kosovo 2010 as \texttt{NaN} since neither existed as a recognized state, and applying linear interpolation for countries with limited macroeconomic gaps of one to four years:
\begin{equation}
x_t = x_{t-1} + \frac{x_{t+k} - x_{t-1}}{k}\cdot i,
\qquad i = 1, \ldots, k-1.
\end{equation}

Three derived variables were constructed. A logarithmic transformation reduced right skew for map color encoding:
\begin{equation}
d_{\log}(c,t) = \log\bigl(\text{total\_displaced}(c,t) + 1\bigr).
\end{equation}

A dynamic binary flag identified the top 30 displacement countries per year:
\begin{equation}
\text{is\_top30}(c,t) =
\begin{cases}
1 & \text{if } \text{rank}_t(c) \leq 30,\\
0 & \text{otherwise.}
\end{cases}
\end{equation}

The Crisis Index combined Min--Max normalized displacement and child mortality into a single 0 to 1 score:
\begin{equation}
\text{CrisisIndex}(c,t) = \frac{1}{2}\left(
\frac{d(c,t) - d_{\min}}{d_{\max} - d_{\min}}
+ \frac{m(c,t) - m_{\min}}{m_{\max} - m_{\min}}
\right).
\end{equation}

Key limitations include: UNHCR data undercounts unregistered populations; UNICEF mortality estimates rely on statistical modelling with wide uncertainty for CAR, DRC, and Somalia; World Bank data is absent for North Korea, Eritrea, and CAR from 2016 onward; the Crisis Index weights displacement and mortality equally, excluding healthcare access and conflict intensity; and the 2010 to 2022 window does not cover the Sudan conflict of April 2023 or the Gaza crisis of October 2023, both of which would alter the rankings substantially.

\section{Audience, Tasks, and Design Rationale}

The system serves two distinct audiences whose needs diverge significantly. The primary audience consists of IFRC communication and advocacy officers who need data-driven evidence to identify priority intervention zones, justify funding allocation, and build evidence-based campaigns. The secondary audience is the general public visiting the campaign webpage---non-expert users who need to understand the scale and geography of the crisis intuitively, without statistical background. This dual audience drives a dual design strategy: analytical depth for the primary user, emotional accessibility for the secondary.

The overall system is organized as a four-part narrative progression hosted on GitHub Pages: global overview (Section~01), crisis identification (Section~02), human stories (Section~03), and ethical reflection (Section~04). Color consistency is maintained across all views, with red encoding severity and grey encoding absence of crisis, reinforcing IFRC's visual identity throughout.

The first two visualizations form a shared Tableau dashboard (T1, T2). A choropleth map colored by log-transformed total displaced was chosen over bubble maps or proportional symbol maps because it communicates geographic patterns at a glance without requiring area estimation. A logarithmic scale was necessary given the extreme right skew of displacement data: a linear scale would render all countries except Syria and Colombia visually identical. A binary choropleth alongside it shifts the user's question from ``how much'' to ``who'', with the top 30 countries recalculated dynamically per year to reflect real shifts in crisis geography. Both maps share a year slider implemented via Tableau's Pages shelf.

The treemap (T4) encodes two quantitative dimensions simultaneously: rectangle size for the Crisis Index and color for child mortality, without requiring axis reading. Restricting it to the 2019 to 2022 average focuses attention on actionable recent data. The design deliberately places CAR and Syria in visual tension, with the largest rectangle belonging to Syria and the darkest color to CAR, communicating that displacement volume and mortality rate are distinct crisis dimensions requiring separate attention.

The bump chart (T2, T5) was chosen over a standard line chart because relative position is more actionable for campaign targeting than absolute value. Filtering to the top 10 countries preserves readability. Ukraine's trajectory from outside the top 10 to third place in 2022 is only legible in a ranking view and would be invisible in a value-based time series. The dot plot (T4) complements the bump chart by providing a scannable absolute ranking colored by geographic region, avoiding the use of the group's one permitted bar chart on a view where it adds limited analytical value.

The pictogram (T1) satisfies the course requirement for an explanatory view accessible to a non-specialist audience. Child icons make mortality rates immediately human and emotionally legible without numerical literacy. The radial timeline satisfies the creative visualization requirement, encoding country identity, year, child mortality, and displacement intensity simultaneously in a polar layout that reinforces the cyclical, recurring nature of humanitarian crises.

The white hat and black hat visualizations use the same CAR child mortality dataset (2015 to 2022) to make the methodological contrast explicit. The white hat version presents the full timeline with a zero-based axis and neutral framing. The black hat version applies four deliberate manipulations: cherry-picked years that erase a key 2019 spike, quadratic smoothing that removes volatility, a truncated y-axis that amplifies the visual slope, and an alarming title that primes emotional response before the reader processes any number.

The JavaScript Country Story Card (T3) bridges the gap between aggregate patterns visible in Tableau and individual country narratives essential for the campaign's call to action. Clicking a country on a Plotly choropleth opens an animated card combining a Chart.js mortality trend line chart (2010 to 2022), a displacement peak bar proportional to the global maximum, a GDP per capita badge, a crisis level badge, and an auto-generated insight sentence. Built in vanilla JavaScript with Plotly and Chart.js loaded from CDN, it is integrated as a custom HTML block in the webpage and satisfies the course requirement for an original JavaScript-based interactive feature.

\section{System Description}

The system is a single-page campaign website hosted on GitHub Pages, structured as a linear narrative across four sections. All visualizations are embedded via \texttt{iframe} directly in the HTML page. The system uses three tools: Tableau Public for five interactive views, Python (\texttt{matplotlib} and \texttt{Plotly}) for static and interactive views, and vanilla JavaScript for the original interactive widget.

\paragraph{Section 01 --- JavaScript Country Story Card.}
The entry point of the system is a Plotly choropleth map coloring each country by crisis level (low, medium, high, critical) derived from the Crisis Index. Clicking a country triggers an animated card containing a Chart.js line chart of child mortality (ages 1 to 4) from 2010 to 2022, a displacement peak bar proportional to the global maximum, a GDP per capita badge, a crisis level badge, and an auto-generated insight sentence computed from the data. Data is loaded from a static \texttt{country\_data.json} file generated by \texttt{export\_json.py} from the master dataset. The widget is built in vanilla JavaScript with Plotly and Chart.js loaded from CDN and requires no backend. This component constitutes the original JavaScript-based interactive feature required by the course.

\paragraph{Section 02 --- Tableau visualizations.}
This section contains three separate Tableau embeds. The first is a dashboard combining two linked choropleth maps built by Chiara Valentini: \textit{Map Displacement}, colored by log-transformed total displaced on a continuous red scale, and \textit{Map Top30 Crisis}, a binary view highlighting the 30 highest-displacement countries per year in red against a grey baseline. Both maps share a Pages shelf year slider covering 2010 to 2022 and are linked so that selecting a year updates both simultaneously. The second embed contains two independent sheets: the bump chart, tracking the Crisis Index ranking of the top 10 countries from 2019 to 2022 with country labels on each line, and the dot plot, ranking all 30 crisis countries by average Crisis Index colored by geographic region. The third embed contains the treemap sized by average Crisis Index and colored by average child mortality for the 2019 to 2022 period. All Tableau views support hover tooltips.

\paragraph{Section 03 --- Python visualizations.}
The explanatory pictogram compares child mortality (ages 1 to 4) between countries in the top and bottom terciles of displacement over 2020 to 2022. Using 50 stylized child icons per row, each colored icon represents one death per 1{,}000 children, making the difference between HIGH displacement countries (13.1 deaths per 1{,}000) and LOW displacement countries (2.1 deaths per 1{,}000) immediately legible without statistical knowledge. The radial timeline displays a polar bar chart for the 10 highest-displacement countries across 2010 to 2022. The chart is organized as a polar wheel with 10 angular slices, one per country, each containing 13 radial bars read counter-clockwise from 2010 to 2022. Bar height encodes child mortality rate and bar color encodes displacement intensity on a logarithmic scale ranging from green to dark red, with the 95th percentile used as the radial maximum to prevent outlier compression. Both views are rendered as self-contained HTML files embedded via \texttt{iframe}.

\paragraph{Section 04 --- White hat and Black hat.}
Two static Python charts are placed side by side using the same dataset: child mortality rates (ages 1 to 4) in the Central African Republic from 2015 to 2022. The white hat chart presents the full timeline with a zero-based axis, all data points including the 2019 spike and 2020 dip, subdued grey labels, and a source attribution. The black hat chart applies four deliberate manipulations: only four cherry-picked years are plotted (2016, 2018, 2020, 2022), erasing the 2019 spike; quadratic spline smoothing connects the selected points, amplifying the visual sweep; the y-axis is truncated at 20 instead of 0; and an alarming title primes emotional response before the reader processes any number. The result transforms a complex, volatile trend into a false narrative of inevitable catastrophe.

The full source code, Python notebooks, JSON data files, and Tableau workbook links are available in the project GitHub repository. A README file documents how to run the Python notebooks, regenerate the JSON export, and access the live Tableau embeds.

\section{Rejected Alternatives}

Two visualization approaches were considered and deliberately rejected. The first was a bubble scatter plot mapping log-transformed displacement on the horizontal axis, child mortality on the vertical axis, GDP per capita as bubble size, and geographic region as color, producing a Gapminder-style view across all 195 countries. When applied to the actual data the design failed on two grounds: the extreme right skew of displacement compressed the majority of countries into an illegible cluster in the low-displacement, low-mortality quadrant, and bubble size encoding for GDP introduced a third perceptual task on top of two already demanding positional judgments. The treemap was adopted in its place, encoding the same dimensions through size and color in a form immediately readable without axis interpretation.

The second rejected alternative was a network graph of refugee flows, with countries as nodes and bilateral displacement flows as directed edges. While the UNHCR dataset's bilateral structure would have supported this design, in practice the graph was visually illegible: with over 190 countries and thousands of flow pairs, nodes overlapped and edges formed an indistinct mass. More fundamentally, the network graph would have reframed the research question from ``where are children dying'' to ``where are people going'', shifting the narrative focus from child health outcomes to migration geography---incompatible with the campaign's objective of identifying countries most in need of pediatric hospital intervention.

\section{Findings and Limitations}

The analysis confirms a systematic but non-linear relationship between forced displacement and child mortality across 2010 to 2022. The data reveals two structurally distinct crisis profiles. The first is displacement-dominant: Syria presents the highest displacement volume across virtually all years yet maintains comparatively moderate child mortality, consistent with a crisis where people flee to countries with functioning healthcare infrastructure. The second is mortality-dominant: the Central African Republic shows relatively moderate displacement but the highest child mortality in the dataset, exceeding 100 deaths per 1{,}000 children aged 5 to 14 in the 2019 to 2022 average. A 2023 nationwide survey estimated that 5.6 percent of CAR's population died in 2022 alone, four times higher than official UN estimates (Karume et al., 2023), demonstrating that displacement volume alone is a poor proxy for child health emergency.

Sub-Saharan African countries cluster in the mid-range of the Crisis Index with structurally high mortality rates that persist independently of displacement levels, suggesting healthcare infrastructure collapse and poverty as confounding variables. GDP per capita correlates with child mortality but not with displacement, confirming that economic development mediates the health impact of forced migration. Ukraine represents the most extreme acute shock in the dataset: absent from the top 10 between 2019 and 2021, it reached third place in 2022 following the Russian invasion---the largest single-year rank jump in the full period. Unemployment and inflation show weaker associations with child mortality than GDP per capita, suggesting structural wealth level matters more than short-term economic volatility.

Several limitations constrain these findings. The Crisis Index weights displacement and mortality equally, a simplification not empirically validated. UNHCR data undercounts undocumented populations. UNICEF mortality estimates for low-income countries rely on statistical modelling with wide confidence intervals---precisely for CAR, DRC, and Somalia. World Bank macroeconomic data is absent for several high-crisis countries across multiple years. The dataset covers 2010 to 2022 and does not include the Sudan conflict escalation of April 2023 or the Gaza crisis of October 2023, both of which would substantially alter the rankings; extending coverage would require sourcing updated UNHCR and UNICEF releases not yet available in comparable format at the time of this analysis.

The white hat and black hat exercise, applied to CAR child mortality from 2015 to 2022, demonstrated that the same dataset can support diametrically opposite narratives. The white hat chart revealed a volatile, non-monotonic trend with a 2019 spike and partial 2020 recovery before a 2022 surge. The black hat version applied four deliberate manipulations---cherry-picked years, quadratic smoothing, a truncated y-axis, and an alarming title---transforming a complex worsening trend into a false narrative of inevitable catastrophe. The implication is direct: design choices in humanitarian data communication are ethical commitments, not aesthetic preferences.

Future improvements would include integrating 2023 data when available, replacing equal weighting with regression-based coefficients, adding healthcare access indicators as a third crisis dimension, and extending the JavaScript Country Story Card with time-series extrapolation.

\newpage
\nocite{*}
\bibliography{biblio.bib}


%%%%%%%%%%%%%%%%%%%%%%%%%%%%%%%%%%%%%%%%%%%%%%%%%%%%%%%%%%%%
\newpage
\appendix

\section{AI usage statement}

This section provides a detailed disclosure of the role Artificial Intelligence (AI) played in the development of this project. AI tools were utilized exclusively as \textbf{assistive technologies} to optimize technical workflows and enhance the clarity of project documentation. The core intellectual property, including the formulation of research questions, the selection and cleaning of datasets, the choice of visual metaphors, and the final analytical and visual interpretations, remains entirely the original work of the authors.

The development process leveraged the \textit{Claude 3.5 Sonnet}, \textit{Claude 3 Opus}, and \textit{Gemini 1.5} models to facilitate the technical execution of the project. Specifically, these models were used to assist in the refinement of complex visualizations across Python, CSS, and JavaScript. This involved adjusting coordinate calculations, debugging data-binding errors, and optimizing code logic to ensure fluid interactivity and responsive design across various display resolutions.

Furthermore, AI was used to streamline the presentation of technical materials by formatting and organizing Markdown content within Jupyter Notebooks and repository documentation. This ensured professional syntax consistency and improved the narrative flow of the technical guides. \textbf{Every AI-generated suggestion or code snippet was subject to rigorous manual review and independent testing to verify accuracy and ensure that the data represented in the visualizations remained untampered.}

No content was generated by AI that would constitute false analytical claims or data hallucinations. Final decisions regarding the implementation of all code and the structure of all documentation were made solely by the authors.


\begin{tcolorbox}[colback=blue!5!white,colframe=blue!75!black,title={Prompt 1: CSS Dashboard Styling (Claude Sonnet)}]
\textbf{Summary:} "Write a CSS stylesheet for the HTML visualization dashboard to ensure a responsive layout with a fixed sidebar and centered chart container." \\
\textbf{Action:} \textit{Accepted.} The AI provided the Flexbox logic and media queries. We implemented the code and manually adjusted the hex codes and spacing to align with the project's visual brand and accessibility standards.
\end{tcolorbox}

\begin{tcolorbox}[colback=blue!5!white,colframe=blue!75!black,title={Prompt 2: LaTeX Report Structure (Claude Sonnet)}]
\textbf{Summary:} "Format this section of the project report in LaTeX, including a professional table for data metrics and properly cited figure captions." \\
\textbf{Action:} \textit{Accepted.} The model converted my raw analysis into clean LaTeX syntax. We used this to ensure consistent document styling and to manage complex table environments that are difficult to format manually.
\end{tcolorbox}

\begin{tcolorbox}[colback=red!5!white,colframe=red!75!black,title={Example of AI Failure \& Correction}]
\textbf{Summary:} "Suggest a method to manipulate the mapping for these black and white visualizations to represent the specific intensity variables." \\
\textbf{Weakness:} The AI output was \textbf{wrong and inconsistent}. Instead of simply adjusting the visual mapping, the AI \textbf{changed the underlying variables} We had specified, replacing my custom data parameters with generic ones. This resulted in a black and white visualization that represented the wrong data entirely. \\
\textbf{Correction:} We \textbf{rejected} the output because it compromised the integrity of our data. We manually rewrote the mapping function to ensure that our original variables were preserved and that the visual hierarchy correctly reflected our specific dataset.
\end{tcolorbox}

%%%%%%%%%%%%%%%%%%%%%%%%%%%%%%%%%%%%%%%%%%%%%%%%%%%%%%%%%%%%




\end{document}